\section{Conclusions and Limitations}
\label{sec:limitations}
We presented SEAL, which learns the localization function of a designed task automaton from a few safe demonstrations, letting multi-step LfD policies recover from task-level perturbations without hand-engineered sensor models or unsafe failure data. The framework rests on a simple observation: with a universal-kernel classifier, model capacity is never the bottleneck: \emph{coverage} is, and the classifier's own epistemic uncertainty marks exactly where coverage is missing, so every expert query converts an uncovered state into a covered one at zero safety cost. By framing the problem as mode localization, SEAL inherits goal reachability from goal-directed primitives and repurposes its calibrated classifier as a data-driven invariance monitor; empirically it matches perturbation-based grounding and expert-tuned automata in accuracy while eliminating their false-positive/latency trade-off, sustaining zero spurious transitions at dwell times where a tuned automaton degrades.

\textbf{Limitations.} (i) High-dimensional observations (images) cannot enter the active loop directly and must be processed into low-dimensional features; recent VLM-based segmentation of human videos~\cite{wang2024vlmseerobotdo} could relax this and enable learning from pre-existing demonstrations. (ii) In the user study, participants could not stage query scenes themselves and needed to see them physically before labeling, so the authors set up each scene, a genuine confound for the effort/preference measures (though not for diversity or mIoU) and a teaching bottleneck; visualizing query points for users is a clear next step. (iii) Performance depends on label quality (cf.\ P09); because we have runtime epistemic uncertainty, calibrating an uncertainty threshold to halt and re-query online is a direct remedy. (iv) We currently assume fixed primitives: if the environment changes so that a subtask always fails, SEAL cannot repair the primitive. The natural extension is to co-train the GPC with imitation-learned policies, which would let us \emph{establish}, rather than assume, mode invariance, and, since active queries target low-density regions, is expected to yield primitives with greater out-of-distribution robustness. (v) The user study is small ($N{=}10$); its head-to-head preference and per-participant accuracy trends are consistent but, as reported in Sec.~\ref{sec:user_study}, not all effects reach significance, and the scene-setup confound limits the effort/preference conclusions.
